\documentclass[preprint,12pt]{elsarticle}

\usepackage{amsmath}
\usepackage{xurl}

% Local paragraph tolerance for long technical terminology.
\setlength{\emergencystretch}{3em}

\journal{Simulation Modelling Practice and Theory}

\begin{document}

\begin{frontmatter}

\title{Reproducible Multi-Layer Evaluation of Cross-Shard Commit Protocols}

\author{Author Name}

\begin{abstract}
Cross-shard transactions require consistent terminal outcomes despite duplicated receipts, quorum failures, timeouts, and partially applied commits. This paper presents a reproducible multi-layer evaluation workflow for an executable cross-shard commit protocol implemented in DTL-Lab. The workflow combines bounded property verification with complementary TLA+ and Alloy models, targeted mutation-based property validation, bounded implementation-model trace conformance, verification-cost characterization, and reproduction of the analytical artifact under a frozen experimental protocol. The definitive campaign scheduled 1,272 tasks. For bounded property verification, 420 measured executions were scheduled, 350 completed, and 70 large-profile TLA+ model-checking executions were censored by timeout; no property violations were observed among the completed configurations. All 10 predefined scientific mutants were detected, yielding a mutation score of 1.0 for the evaluated catalogue. Across 600 trace-conformance executions, the valid population was accepted and the deliberately corrupted population was rejected with the expected diagnostics. Alloy median elapsed time increased from 0.515306 to 0.766172 and 1.568320 seconds across the three bound profiles, while all measured large-profile TLA+ executions reached the timeout threshold. A separate clone and workspace on the same native Linux host completed 10/10 reproduction gates and matched 32/32 expected analysis-artifact hashes. The results provide complementary bounded evidence and an empirical verification-cost boundary; they do not constitute an unbounded correctness proof, general refinement, tool-superiority claim, nonlinear scaling law, or hardware-independent reproduction.
\end{abstract}

\begin{keyword}
cross-shard transactions \sep
distributed ledgers \sep
formal verification \sep
model checking \sep
mutation analysis \sep
trace conformance \sep
reproducibility
\end{keyword}

\end{frontmatter}

\section{Introduction}
\label{sec:introduction}

Sharding increases the capacity of distributed ledgers by partitioning state
and transaction processing across multiple shards. This parallelism creates a
coordination problem whenever a transaction spans more than one shard: the
system must preserve a consistent terminal outcome despite delayed or
duplicated messages, unmet quorum conditions, or interrupted execution.
Chainspace, OmniLedger, Prophet, CSLAP, and LightCross illustrate different
approaches to cross-shard transaction processing and atomic commit
\cite{albassam2018chainspace,kokoriskogias2018omniledger,hong2023prophet,du2024cslap,qi2024lightcross}.
Their correctness arguments, implementation evidence, and evaluation
methodologies differ substantially.

Formal methods provide evidence about protocol behavior beyond conventional
testing. TLA+ and related model-checking workflows analyze distributed
protocols and connect formal models with executable systems
\cite{schultz2025mongodb,cirstea2025traces,hackett2025tracelink,howard2025ccf,wang2023mocket,gao2025imocket}.
Mutation-based assessment instead asks whether a specification or property
suite reacts when selected controls are intentionally weakened
\cite{cordy2023mutation}. These techniques are complementary: bounded model
checking evaluates declared properties in finite configurations, mutation
analysis measures sensitivity to controlled defects, and trace conformance
checks selected implementation observations against a formal abstraction.

This paper studies an executable cross-shard commit protocol implemented in
DTL-Lab, a research-oriented Java simulator rather than a production
blockchain. The protocol coordinates source locking, receipt creation and
delivery, destination preparation, commit, abort, timeout, rollback, replay
protection, and quorum enforcement. Bounded TLA+ and Alloy representations
model the same protocol concerns through different state representations and
exploration semantics. The study assumes neither semantic equivalence between
the formal models nor intrinsic superiority of either tool.

The methodological objective is to connect several forms of evidence under one
frozen experimental protocol. Within the comparator set for which full text was
verified in our structured review, we did not identify a cross-shard study
reporting together bounded property verification with complementary TLA+ and
Alloy models, targeted mutation-based property validation, bounded
implementation-model trace conformance over valid and deliberately corrupted
traces, explicit verification-cost characterization, and independent
reproduction of the analytical artifact. This statement is limited to the
reviewed comparator set and is not a priority claim about the complete
literature. The contribution is therefore an integrated and reproducible
evidence workflow rather than novelty of any individual technique.

The study is organized around four research questions:

\begin{description}
\item[\textbf{RQ1.}] What bounded property-verification outcomes are obtained
for the valid TLA+ and Alloy models across the declared experimental
configurations?

\item[\textbf{RQ2.}] Do the evaluated TLA+ and Alloy properties detect
predefined scientific mutations that remove selected protocol controls?

\item[\textbf{RQ3.}] Do selected valid and deliberately corrupted Java
executions satisfy the expected bounded implementation-model trace-conformance
outcomes under the declared projection and scenario catalogue?

\item[\textbf{RQ4.}] How does verification cost vary across the evaluated
bound profiles and TLC fault profiles under the frozen execution protocol?
\end{description}

The evidence is intentionally layered. RQ1 provides bounded property
verification, RQ2 evaluates sensitivity to predefined scientific mutants, RQ3
evaluates bounded implementation-model trace conformance without claiming
refinement or behavioral equivalence, and RQ4 characterizes verification cost
within each tool without implying an asymptotic law or absolute
TLC-versus-Alloy superiority.

The definitive campaign scheduled 1,272 tasks. RQ1 included 420 measured
executions, of which 350 completed and 70 TLC-large executions were censored by
timeout, with no property violations observed among completed configurations.
RQ2 detected all 10 predefined scientific mutants, yielding a mutation score of
1.0 for the evaluated catalogue. RQ3 comprised 600 executions covering valid
and deliberately corrupted traces, valid traces were accepted and corrupted
traces rejected with the expected diagnostics. RQ4 showed increasing
verification cost across completed Alloy profiles, while all measured
TLC-large executions reached timeout. These observations support bounded and
tool-specific conclusions only.

Reproduction used a separate clone, execution workspace, and operating-system
user on the same native Linux host. The successful attempt completed 10/10
predefined gates and reproduced 32/32 expected analysis-artifact hashes. This
demonstrates process separation and deterministic reconstruction of preserved
analytical outputs under the documented environment, not hardware independence,
a complete rerun of all 1,272 tasks, or scientific correctness from hash
agreement alone.

The remainder of the paper presents related work, the protocol and formal
models, methodology and experimental design, results and interpretation,
threats to validity, and reproducibility evidence.


\section{Background and Related Work}
\label{sec:background-related-work}

The literature relevant to this study spans cross-shard transaction protocols,
formal analysis of distributed and blockchain systems, implementation-model
validation, and reproducibility of verification artifacts. These areas contain
mature techniques used individually in our workflow, so we position the
contribution in their integration and controlled evaluation for one executable
cross-shard protocol.

\subsection{Sharded ledgers and cross-shard commit}
\label{subsec:related-cross-shard}

Sharded ledgers have long treated atomic cross-shard processing as a central
systems problem. Chainspace introduced S-BAC, a Byzantine distributed atomic
commit protocol for transactions whose objects span multiple shards, together
with correctness arguments and an executable prototype
evaluation~\cite{albassam2018chainspace}. OmniLedger introduced Atomix, a
client-driven atomic cross-shard transaction protocol that coordinates input
and output shards using proofs of acceptance or
rejection~\cite{kokoriskogias2018omniledger}. These systems establish that
atomic cross-shard commit mechanisms and protocol-level correctness arguments
predate the present work.

More recent systems focus on reducing conflicts, latency, or coordination
overhead. Prophet uses Byzantine-tolerant deterministic ordering to reduce
cross-shard conflicts and aborts and provides arguments for determinism and
serializability~\cite{hong2023prophet}. CSLAP builds leader accountability on
top of two-phase atomic commit and analyzes consistency, liveness, and
communication behavior~\cite{du2024cslap}. LightCross combines TEE-assisted
execution with a lightweight cross-shard commit protocol and presents
atomicity, serializability, and liveness arguments for its
design~\cite{qi2024lightcross}.

The cited cross-shard systems already provide correctness evidence through
protocol design, analytical or theorem-based arguments, and system-level
evaluation using metrics such as throughput and latency. Those metrics
characterize implementation behavior and should not be conflated with the
resource cost of formal verification. In contrast, RQ4 characterizes bounded
model-checking cost under a predefined experimental protocol. Mutation-based
property validation and implementation-model trace conformance are also not
reported as evaluation dimensions in the cited cross-shard works.

\subsection{Formal analysis of distributed and blockchain protocols}
\label{subsec:related-formal-analysis}

Formal modeling and model checking have been applied directly to distributed
transaction protocols. Schultz and Demirbas model MongoDB distributed
transactions in TLA+ and use TLC to analyze protocol properties. Their modular
treatment additionally generates tests from a formal storage model and
executes them against WiredTiger, providing a direct connection between the
model and a production implementation interface~\cite{schultz2025mongodb}.
This work is particularly close to our study because it combines cross-shard
distributed transactions, bounded formal exploration, and implementation-facing
validation.

Mutation-based assessment of formal specifications also predates this work.
Cordy et al. propose mutation model checking to assess and strengthen
specifications by checking whether deliberately mutated models are rejected by
the property set~\cite{cordy2023mutation}. Their work establishes
mutation-based specification assessment as a prior methodological approach.
Our RQ2 uses a smaller, protocol-targeted catalogue of scientific mutants and
interprets its mutation score only as sensitivity to those predefined defect
classes.

The use of multiple formal encodings and deliberate defects is likewise not
unique to this study. Konnov et al. translate the executable Ethereum 3SF
specification into TLA+, analyze it with Apalache, and use alternative SMT and
Alloy encodings for cross-validation~\cite{konnov2025ethereum}. Their technical
report also evaluates time, memory, timeouts, and deliberately injected bugs.
Consequently, neither the combination of TLA+ and Alloy nor the combination of
multi-formalism verification, defect injection, and verification-cost
measurements should be treated as an isolated novelty claim in our manuscript.

\subsection{Specification validation and implementation-model conformance}
\label{subsec:related-conformance}

A substantial body of work addresses the gap between verified models and
executable implementations. Cirstea et al. reduce validation of
distributed-program traces against TLA+ specifications to constrained model
checking with TLC and provide an API for instrumenting Java
programs~\cite{cirstea2025traces}. Their framework permits partially observed
traces and makes explicit the tradeoff between instrumentation detail and the
search required to determine whether an observed execution is compatible with
the specification.

TraceLink further automates the relation between distributed implementations
and verified designs. It links instrumented executions to TLA+ models,
evaluates trace-validation strategies, reports previously undetected
implementation defects, and publishes an experimental artifact with partial
external reproduction~\cite{hackett2025tracelink}. Smart Casual Verification
of CCF applies a related approach to a production C++ codebase by integrating
TLA+ model checking, simulation, testing, and trace validation into continuous
integration~\cite{howard2025ccf}. These studies show implementation-model trace
validation across multiple distributed-systems settings.

Other approaches generate or guide implementation checks from formal models.
Choreographic PlusCal projects distributed protocol descriptions to TLA+ and
generates conformance monitors, including evaluation on two-phase
commit~\cite{foo2023choreographic}. Mocket uses a verified TLA+ state space to
generate systematic tests, drive Java implementations through selected
behaviors, and compare runtime states with formal
counterparts~\cite{wang2023mocket}. iMocket restricts this strategy to
state-space regions affected by incremental changes~\cite{gao2025imocket}.
Beyond TLA+, the Amazon S3 ShardStore study combines executable reference
models, property-based testing, failure injection, and stateless model
checking~\cite{bornholt2021s3}.

These results constrain RQ3. We do not present trace validation or
implementation-model checking as a new general method. RQ3 evaluates bounded
implementation-model trace conformance for a predefined cross-shard catalogue
containing valid and deliberately corrupted traces with expected diagnostic
agreement. Its role is one implementation-facing layer in the evidence chain
used for RQ1, RQ2, and RQ4.

\subsection{Positioning of this work}
\label{subsec:related-positioning}

Within the comparator set for which full text was verified in our structured
review, we did not identify a cross-shard study that reports all of the
following within one frozen experimental protocol: bounded property
verification using complementary TLA+ and Alloy models, targeted mutation-based
property validation, bounded implementation-model trace conformance over both
valid and deliberately corrupted traces, explicit verification-cost
characterization, and independent reproduction of the analytical artifact.
This statement is deliberately scoped to the reviewed comparator set and is
not a priority claim about the complete literature.

Accordingly, DTL-Lab is positioned as an integrated and reproducible evidence
workflow for an executable cross-shard protocol. It connects bounded formal
verification, mutation sensitivity, implementation-model conformance, and
verification cost while preserving each layer's limitations, without claiming
the individual techniques as novel.


\section{Cross-Shard Transaction Model}
\label{sec:cross-shard-model}

\subsection{System abstraction}
\label{subsec:system-abstraction}

We study a cross-shard transaction protocol in which a transfer moves value from
a source shard to a target shard while preserving a consistent terminal outcome.
Each shard maintains an independent UTXO-based ledger, and each cross-shard
transfer is represented by a uniquely identified protocol session.

A session coordinates source locking, cross-shard receipt creation and delivery,
destination preparation, and the final commit decision, while also handling
aborts, validation failures, insufficient quorum, timeouts, duplicate receipts,
and partially applied commits.

The implementation separates shard management from transaction-state mutation.
The shard manager maintains topology, validators, logical time, session
registration, and aggregate metrics, while the atomic-commit component validates
protocol preconditions and applies state changes. This separation allows the
transaction protocol to be analyzed independently from unrelated functionality
in the surrounding simulation platform.

Receipt authenticity is an abstract protocol precondition rather than a
cryptographic result. Quorum summarizes validator approval without modeling the
consensus mechanism that produces votes, and timeouts use logical rounds rather
than wall-clock deadlines.

\subsection{Cross-shard session state machine}
\label{subsec:state-machine}

Each transfer is represented by a finite-state session. The implementation
contains the following states:

\begin{itemize}
\item \texttt{CREATED}
\item \texttt{SOURCE\_\allowbreak LOCKED}
\item \texttt{RECEIPT\_\allowbreak CREATED}
\item \texttt{RECEIPT\_\allowbreak DELIVERED}
\item \texttt{DESTINATION\_\allowbreak PREPARED}
\item \texttt{COMMITTED}
\item \texttt{ABORTED}
\item \texttt{TIMED\_\allowbreak OUT}
\item \texttt{FAILED\_\allowbreak VALIDATION}
\end{itemize}

The normal successful execution follows

\begin{equation}
\begin{split}
\texttt{CREATED}
&\rightarrow \texttt{SOURCE\_\allowbreak LOCKED} \\
&\rightarrow \texttt{RECEIPT\_\allowbreak CREATED} \\
&\rightarrow \texttt{RECEIPT\_\allowbreak DELIVERED} \\
&\rightarrow \texttt{DESTINATION\_\allowbreak PREPARED} \\
&\rightarrow \texttt{COMMITTED}.
\end{split}
\label{eq:normal-commit-path}
\end{equation}

The states \texttt{COMMITTED}, \texttt{ABORTED}, \texttt{TIMED\_\allowbreak OUT},
and \texttt{FAILED\_\allowbreak VALIDATION} are terminal and have no outgoing
transitions.

Abort and validation-failure transitions may occur from nonterminal states. A
timeout may occur after source locking and before a terminal decision, but not
directly from \texttt{CREATED}, where no source resource has been locked.

The transition relation rejects commit after abort, commit after timeout,
creation-to-commit jumps, repeated terminal decisions, and receipt processing
for a different transfer identifier. Each accepted transition emits an
immutable event containing sequence number, logical time, transfer identifier,
action, previous state, next state, and diagnostic reason.

\subsection{Atomic commit and rollback}
\label{subsec:atomic-commit}

The executable commit operation has four conceptual stages:

\begin{enumerate}
\item preparation of the commit plan
\item validation of commit preconditions
\item application of ledger and session mutations
\item rollback if an application step fails
\end{enumerate}

Before modifying ledger state, the protocol constructs a commit plan and
captures a restorable snapshot containing the source UTXO and lock state,
receipt-consumption state, optional change output, destination output, and a
checkpoint of the complete protocol session.

During normal commit, the destination is prepared, the receipt consumed, the
source value removed, the source lock released, optional change created, the
destination output created, and the session moved to \texttt{COMMITTED}.

A failure invokes rollback before the protocol error is propagated. Rollback
restores the destination and change outputs, source UTXO, receipt-consumption
state, source lock, and session checkpoint. No synthetic reverse transition is
recorded because rollback cancels an incomplete state application rather than
creating a new protocol decision.

Deterministic failure-injection points around receipt consumption, source debit,
and destination credit make partially executed commits reproducible without
introducing nondeterministic faults.

\subsection{Concurrent formal abstraction}
\label{subsec:formal-abstraction}

Bounded TLA+ and Alloy models abstract the executable protocol to the
protocol-level state required for concurrent sessions and failure behavior. They
represent transfer status, source and target shards, resource locks, receipt
ownership and consumption, destination credit, released funds, messages, and
validator votes, while collapsing implementation states that do not alter the
protocol-level decision state.

The Java-to-formal projection maps \texttt{CREATED} to \texttt{Pending};
\texttt{SOURCE\_\allowbreak LOCKED}, \texttt{RECEIPT\_\allowbreak CREATED}, and
\texttt{RECEIPT\_\allowbreak DELIVERED} to \texttt{Locked};
\texttt{DESTINATION\_\allowbreak PREPARED} to \texttt{Prepared};
\texttt{COMMITTED} to \texttt{Committed}; and \texttt{ABORTED},
\texttt{TIMED\_\allowbreak OUT}, and \texttt{FAILED\_\allowbreak VALIDATION} to
\texttt{Aborted}.

TLA+ additionally stores the first abstract terminal decision in
\texttt{terminalStatus}. Alloy instead represents bounded executions as an
ordered sequence of states and checks terminal-state irreversibility over
consecutive states. The two models therefore encode terminal consistency
differently while targeting the same protocol concern.

Concurrent configurations vary shards and active transfers and include
duplicate receipts, delayed messages, timeout races, and quorum behavior. They
are finite by construction, so conclusions are bounded by the declared scopes
and constants.

TLA+ and Alloy are complementary representations of the same protocol concerns,
not equivalent state spaces or interchangeable cost models. Absolute
execution-time differences are therefore not used to infer tool superiority.

\subsection{Evaluated protocol properties}
\label{subsec:properties}

Seven protocol properties define the main bounded verification target.

\begin{description}

\item[] \texttt{NoReceipt\allowbreak Replay}.
A cross-shard receipt must not be consumed more than once.

\item[] \texttt{DestinationCreditRequires\allowbreak ValidReceipt}.
Destination-side credit requires a receipt accepted as valid by the protocol.

\item[] \texttt{Decision\allowbreak Consistency}.
A committed transfer must not simultaneously be represented as having released
its source funds through an abort-like terminal outcome.

\item[] \texttt{NoValueLossAt\allowbreak Termination}.
At a terminal state, transferred value must either have been committed to the
destination or remain recoverable on the source side.

\item[] \texttt{TerminalState\allowbreak Irreversibility}.
Once the protocol reaches its first terminal decision, a later execution step
must not replace it with an incompatible terminal outcome.

\item[] \texttt{EventuallyReleasedAfter\allowbreak Timeout}.
An abstract transfer represented as \texttt{Aborted} must have released its
source funds. This property captures the release obligation associated with the
modeled timeout and non-commit terminal behavior.

\item[] \texttt{Quorum\allowbreak Required}.
A commit requires the relevant validator quorum. TLA+ parameterizes this
threshold through \texttt{Quorum}, whereas the evaluated Alloy model requires at
least two validator votes.

\end{description}

Together these properties cover replay protection, destination authorization,
decision consistency, value conservation, terminal-state stability, timeout
release, and quorum enforcement. The same set is evaluated for valid formal
models and the mutation-based experiments.

\subsection{Scope of the model}
\label{subsec:model-scope}

Receipt authenticity remains abstract, validator quorum is an approval
condition rather than a verified consensus protocol, and time is represented by
logical rounds. Network behavior is scenario-bounded: delay and duplicate
receipts are represented in evaluated configurations, while concrete network
events not structurally represented by the abstraction are projected as
stuttering steps during trace conformance.

Both formal representations operate over finite configurations. Absence of a
counterexample therefore means only that no violation was found for the
evaluated model, property, configuration, and bound, not an unbounded proof of
protocol correctness.

The Java implementation is related to the formal model through bounded trace
conformance, not a general refinement proof or behavioral equivalence.


\section{Research Methodology}
\label{sec:methodology}

\subsection{Study design}
\label{subsec:study-design}

This study evaluates a cross-shard commit protocol through bounded property
verification, mutation-based property validation, bounded implementation-model
trace conformance, and verification-cost characterization. The four methods
address complementary questions rather than serving as interchangeable
evidence.

The design was frozen before the definitive experimental campaign. Research
questions, hypotheses, factors, configuration profiles, seeds, repetitions,
resource limits, statistical procedures, and rules for incomplete executions
were specified before observing the final results. Changes to these elements
after inspection of the definitive outcomes would require a new protocol
version.

\subsection{Research questions and hypotheses}
\label{subsec:research-questions}

\paragraph{RQ1: Preservation of declared properties.}

Does the cross-shard commit protocol preserve the declared properties under the
explored interleavings within the documented bounds?

The corresponding hypothesis is:

\begin{quote}
\textbf{H1.} Valid models produce no property violations within the declared
bounds.
\end{quote}

RQ1 is interpreted only for the executed property, tool, configuration, and
bound, not as an unbounded correctness proof.

\paragraph{RQ2: Defect-detection capability.}

Do TLC and Alloy detect predefined scientific mutations that remove controls
from the cross-shard protocol?

The corresponding hypothesis is:

\begin{quote}
\textbf{H2.} Each scientific mutant produces a violation of its designated
target property.
\end{quote}

Detection requires the designated target property to be reported among the
violations, although a mutant may violate additional properties.

\paragraph{RQ3: Bounded implementation-model trace conformance.}

Do observable traces produced by the Java implementation correspond to action
sequences admitted by the TLA+ model within the evaluated scenarios and bounds?

The corresponding hypothesis is:

\begin{quote}
\textbf{H3.} Valid traces are accepted and corrupted traces are rejected at the
expected point.
\end{quote}

RQ3 provides bounded implementation-model trace-conformance evidence rather
than a general refinement proof or behavioral equivalence.

\paragraph{RQ4: Verification cost.}

How does verification cost change as the number of shards, concurrent
transfers, bounds, and enabled faults increase?

The preregistered hypothesis is:

\begin{quote}
\textbf{H4.} Model-checking cost grows nonlinearly with the number of concurrent
transfers and shards.
\end{quote}

H4 is empirical rather than assumed and is evaluated from observed
profile-level cost measurements.

\subsection{Methodological layers}
\label{subsec:methodological-layers}

The four questions form three verification and validation layers plus one
cost-analysis layer.

\subsubsection{Bounded property-verification layer}

RQ1 evaluates the valid TLA+ and Alloy specifications against the seven
protocol properties in Section~\ref{subsec:properties}. The experimental unit
consists of:

\begin{itemize}
\item a protocol property
\item a formal tool
\item a bounded configuration
\item a measured repetition
\end{itemize}

Response variables include property outcome, available state-space information,
exploration depth or scope, elapsed time, maximum resident memory, exit status,
and counterexample information. A completed execution without a counterexample
provides evidence only within its declared finite configuration. Timeout and
out-of-memory outcomes remain incomplete or censored observations.

\subsubsection{Mutation-based property validation layer}

RQ2 evaluates scientific mutants that remove specific protocol controls:
replay protection, receipt-dependent destination credit, terminal-decision
consistency, timeout release of source funds, and quorum enforcement. The
experimental unit consists of:

\begin{itemize}
\item a scientific mutant
\item a formal tool
\item a bounded configuration
\item a measured repetition
\end{itemize}

Each mutant has one predefined target property, and detection is credited only
when that property is violated. The aggregate measure is

\begin{equation}
\mathrm{MutationScore}
=
\frac{N_{\mathrm{detected}}}
{N_{\mathrm{mutants}}},
\label{eq:mutation-score}
\end{equation}

where \(N_{\mathrm{detected}}\) is the number of predefined mutants detected
through their target properties and \(N_{\mathrm{mutants}}\) is the number
evaluated. The score applies only to the predefined catalogue and does not
establish property-suite completeness.

\subsubsection{Bounded implementation-model trace-conformance layer}

RQ3 evaluates deterministic Java traces projected through the Java-to-TLA+
abstraction in Section~\ref{subsec:formal-abstraction}. The experimental unit
consists of:

\begin{itemize}
\item a valid scenario or negative mutation
\item a deterministic seed
\item the corresponding TLC replay outcome
\end{itemize}

Two trace classes are evaluated: valid traces generated from supported protocol
scenarios and corrupted traces containing predefined structural or semantic
defects. A valid trace is classified successfully when its projected action
sequence is admitted by the replay model. A corrupted trace succeeds
experimentally when it is rejected at the predefined diagnostic point.

For negative traces, the observed diagnostic must also agree with the expected
abstract step, concrete step, action, and transfer identifier. These criteria
distinguish generic replay failure from rejection at the intended protocol
location. The evidence remains bounded to the frozen scenario, mutation, seed,
and projection catalogue.

\subsubsection{Verification-cost characterization layer}

RQ4 characterizes bounded verification cost. The experimental unit consists of:

\begin{itemize}
\item a formal tool
\item a bound profile
\item a fault profile when applicable
\item a measured repetition
\end{itemize}

Response variables are elapsed time, maximum resident memory, generated and
distinct states when available, exploration depth or Alloy scope, timeout,
out-of-memory outcome, and within-tool growth ratios. Because TLC and Alloy
expose different exploration semantics and state representations, cost trends
are analyzed within each tool rather than by treating their absolute
measurements as interchangeable.

\subsection{Triangulation of evidence}
\label{subsec:evidence-triangulation}

The layers address different weaknesses in the scientific argument. RQ1
evaluates bounded property outcomes but cannot show by itself that the
properties expose relevant defects. RQ2 adds targeted defect sensitivity.
Formal-model evidence alone does not establish agreement with the executable
implementation, so RQ3 adds bounded trace conformance with positive and
negative cases. RQ4 characterizes the computational envelope in which the
formal evidence is obtained.

The dependency among the layers is deliberate. Successful bounded checking in
RQ1 is more informative when RQ2 shows that selected weakened controls trigger
their target properties. Those two layers remain statements about formal
representations, so RQ3 confronts the abstraction with executable traces. RQ4
then reports the resource envelope within which the formal evidence was
obtained.

Thus, the first three layers provide complementary evidence about bounded
property outcomes, defect sensitivity, and implementation-model trace
conformance, while the fourth characterizes their verification cost. The layers
strengthen different parts of the argument, but no layer is interpreted as
replacing the limitations of another.

\subsection{Treatment of incomplete executions}
\label{subsec:incomplete-executions}

Experimental outcomes are retained according to observed status.
Counterexamples, timeouts, out-of-memory results, and other scientifically
meaningful unfavorable outcomes are not removed. A failed instrumentation
attempt may be repeated only under the predefined protocol rules, not because a
run produced an unexpected scientific result.

Timeout and out-of-memory observations are retained as censored results rather
than assigned artificial elapsed-time or memory values or reclassified as
successful verification. This prevents incomplete runs from being converted
into positive evidence.

\subsection{Interpretation boundaries}
\label{subsec:method-boundaries}

All claims remain bounded by the declared protocol model, properties,
configurations, scopes, scenarios, seeds, mutant catalogue, and execution
environment. Absence of a counterexample is not an unrestricted proof,
mutation analysis does not establish property-suite completeness, and trace
replay does not establish general refinement or behavioral equivalence. TLC and
Alloy have different semantics and cost units, so their absolute measurements
do not establish intrinsic tool superiority. The evaluated configurations are
abstractions rather than unrestricted production blockchain systems.


\section{Experimental Design}
\label{sec:experimental-design}

\subsection{Frozen experimental protocol}
\label{subsec:frozen-protocol}

The definitive experiment follows protocol \texttt{paper1-q3-v1}, frozen before
execution of the experimental matrix. It specifies the research questions,
hypotheses, configuration families, seeds, repetitions, resource limits,
statistical procedures, and treatment of incomplete executions.

The scientific baseline corresponds to release \texttt{v1.1.0-rc.1}. The frozen
toolchain consists of TLA+ Tools 1.7.4, Alloy 6.2.0 with SAT4J, Java 17,
Python 3.12, and GNU \texttt{/usr/bin/time -v} for time and memory
instrumentation.

The definitive design covers bounded property verification, scientific mutants,
bounded implementation-model trace conformance, and TLC fault profiles.
Configuration families were enumerated before definitive execution. No
configuration absent from the frozen matrix was added after observing outcomes.

\subsection{Configuration matrix}
\label{subsec:configuration-matrix}

The experimental matrix contains fourteen configuration families:

\begin{itemize}
\item six valid-model scalability families, with small, medium, and large
profiles for TLC and Alloy
\item two mutation families, one per formal tool
\item two trace-conformance families, valid and negative
\item four TLC fault-profile families covering normal execution, replay
pressure, timeout behavior, and insufficient quorum
\end{itemize}

These families expand through properties, mutants, repetitions, scenarios, and
seeds into 1272 scheduled tasks. The count includes warm-up and measured
executions where defined, while RQ3 contributes one replay execution per
case-seed pair.

Each task records protocol and configuration identifiers, tool, model kind,
profile, repetition or seed, resource limits, and provenance metadata.

\subsection{Bound profiles}
\label{subsec:bound-profiles}

Three ordered profiles define the principal bounded scalability configurations
in Table~\ref{tab:bound-profiles}.

\begin{table}[htbp]
\centering
\footnotesize
\setlength{\tabcolsep}{2pt}
\caption{Bound profiles used in the definitive experiment.}
\label{tab:bound-profiles}
\begin{tabular}{lrrrrrrr}
\hline
Profile &
Shards &
Transfers &
Validators &
\shortstack{TLC\\quorum} &
\shortstack{TLC receipt\\copies} &
\shortstack{Alloy Receipt\\scope} &
\shortstack{Alloy State\\scope} \\
\hline
Small  & 2 & 1 & 2 & 2 & 1 & 2 & 6  \\
Medium & 3 & 2 & 3 & 2 & 2 & 4 & 8  \\
Large  & 4 & 3 & 4 & 3 & 2 & 6 & 10 \\
\hline
\end{tabular}
\end{table}

Shard, transfer, and validator counts are shared across the profile records.
TLC quorum and receipt-copy columns are TLC parameters. Alloy instead uses
Receipt scopes 2, 4, and 6, Message scopes 10, 20, and 30, and State scopes
6, 8, and 10 for small, medium, and large. Its commit predicate requires at
least two votes in every profile, so TLC quorum values are not Alloy quorum
thresholds.

Profiles are experimental bundles, not one-factor perturbations: shared
configuration sizes and tool-specific bounds vary together. Cost changes are
therefore associated with the complete profile rather than attributed causally
to one parameter. The \texttt{catalog} profile is reserved for RQ3 scenarios.

\subsection{Mutation configurations}
\label{subsec:mutation-configurations}

Mutation experiments use the medium profile. Five scientific defect classes are
instantiated for each formal representation:

\begin{itemize}
\item removal of receipt replay protection
\item destination credit before valid receipt consumption
\item commit after an abort-like terminal decision
\item timeout without release of source funds
\item quorum bypass
\end{itemize}

Each mutant has a predefined target property. After two warm-ups, ten measured
repetitions are executed per mutant, and logical detection must remain
consistent across measured repetitions.

\subsection{Trace-conformance configurations}
\label{subsec:trace-configurations}

RQ3 uses a separate multiseed design with ten valid scenarios and ten negative
cases. Each case is evaluated under thirty deterministic seeds, producing
300 valid and 300 negative trace executions. The definitive seed sequence is
\texttt{2026001} through \texttt{2026030}.

RQ3 therefore contains 600 TLC replay executions. Each case-seed pair executes
once without warm-up because the primary outcome is deterministic
classification, not repeated performance estimation.

\subsection{Repetitions and execution order}
\label{subsec:repetitions}

Model-checking configurations used for time and memory measurement have two
warm-up repetitions, ten measured repetitions, serial execution, and exclusion
of warm-ups from primary statistical analysis. Warm-up observations remain in
the execution record but do not contribute to the reported cost summaries.

Logical outcomes are expected to remain consistent across repetitions.
Repetition characterizes elapsed time and memory rather than creating
independent logical property outcomes. No run is repeated because it produces a
counterexample, timeout, or other unfavorable scientific result.

\subsection{Resource limits}
\label{subsec:resource-limits}

Each measured model-checking execution uses:

\begin{itemize}
\item timeout: 1800 seconds
\item maximum resident-memory limit: 12288 MiB
\item maximum concurrent experimental executions: one
\item TLC workers: one
\item Alloy solver: SAT4J
\end{itemize}

Time and maximum resident memory are collected with
\texttt{/usr/bin/time -v}. Timeout and out-of-memory outcomes remain
scientifically meaningful scalability observations in the experimental record
and are not removed from the matrix.

\subsection{Execution environment}
\label{subsec:execution-environment}

The definitive timing and memory campaign runs on a dedicated native Linux host.
Virtualized environments, WSL, and shared continuous-integration runners are
excluded from primary timing and memory measurements. Continuous integration is
used for functional, structural-integrity, and reproducibility checks, but its
timings are not mixed with definitive performance measurements.

Recorded environment metadata includes operating system and kernel, processor
architecture and count, available memory, Java and Python versions, formal-tool
versions, repository commit, and virtualization status.

\subsection{Measured outcomes}
\label{subsec:measured-outcomes}

For TLC, recorded measurements include:

\begin{itemize}
\item property result
\item generated states
\item distinct states
\item exploration depth
\item elapsed time
\item maximum resident memory
\item process exit status
\item counterexample information when applicable
\end{itemize}

For Alloy, recorded measurements include:

\begin{itemize}
\item assertion result
\item solver and scope
\item counterexample information
\item solving and total elapsed time when available
\item maximum resident memory
\item process exit status
\end{itemize}

TLC state counts and Alloy scopes are not equivalent state-space units.

Trace-conformance records include case identifier, seed, case type, acceptance
or rejection, diagnostic agreement, abstract and concrete steps, action,
transfer identifier, elapsed time, memory, source commit, and input hashes.

\subsection{Incomplete executions and censoring}
\label{subsec:censoring}

Executions may terminate as completed runs, counterexamples, timeouts,
out-of-memory outcomes, tool errors, or instrumentation failures. Timeout and
out-of-memory outcomes are retained as censored observations without replacing
their elapsed-time or memory values by configured limits or synthetic values.

Unexpected tool exits remain tool errors. Instrumentation failures may be
retried only under the frozen retry policy and never on the basis of the
scientific result. This treatment preserves unfavorable outcomes and prevents
selective removal of executions that would weaken a hypothesis.

\subsection{Statistical analysis}
\label{subsec:statistical-analysis}

Time and memory distributions are summarized by median and interquartile range,
with minimum, maximum, and coefficient of variation also retained. A 95\%
bootstrap confidence interval for the median uses 10,000 resamples. Wilson 95\%
confidence intervals are used for proportions, including mutation detection and
trace classification.

RQ4 reports profile-level medians and within-tool growth ratios. Spearman rank
association describes monotonic association between ordered profile and cost
when sufficient completed profiles are available.

TLC and Alloy are not compared through absolute execution-time claims because
their modelling semantics and state-space representations differ. Such
measurements are not treated as evidence of intrinsic tool superiority. The
three ordered profiles provide descriptive bounded-cost evidence but do not
establish an asymptotic complexity class or a general nonlinear growth law.


\section{Results}
\label{sec:results}

\subsection{Experimental campaign overview}
\label{subsec:results-overview}

The definitive campaign contained 1272 scheduled experimental tasks
across the four research questions. This section reports only measured
runs when evaluating the research hypotheses; warm-up executions are
reported separately from the principal statistical evidence.

The results are presented in the order of the research questions.
Incomplete executions are retained according to the censoring policy
defined in Section~\ref{subsec:censoring}.

\subsection{RQ1: Bounded property verification}
\label{subsec:rq1-results}

RQ1 asks whether the valid formal models satisfy the seven declared
protocol properties within the evaluated bounds.

Across both formal tools and the three bound profiles, RQ1 contained
420 measured property runs. Of these, 350 completed and 70 were
censored by timeout. No completed run produced a counterexample.

Table~\ref{tab:rq1-summary} summarizes the result by tool and bound
profile.

\begin{table}[htbp]
\centering
\caption{RQ1 bounded property-verification results by tool and profile.}
\label{tab:rq1-summary}
\begin{tabular}{llrrrr}
\hline
Tool & Profile & Measured & Completed & No counterexample & Timeout \\
\hline
Alloy & Small  & 70 & 70 & 70 & 0  \\
Alloy & Medium & 70 & 70 & 70 & 0  \\
Alloy & Large  & 70 & 70 & 70 & 0  \\
TLC   & Small  & 70 & 70 & 70 & 0  \\
TLC   & Medium & 70 & 70 & 70 & 0  \\
TLC   & Large  & 70 & 0  & 0  & 70 \\
\hline
Total &        & 420 & 350 & 350 & 70 \\
\hline
\end{tabular}
\end{table}

For Alloy, all seven properties completed without a counterexample
under the small, medium, and large profiles. Each property-profile
combination was measured ten times, yielding 210 completed Alloy runs.

For TLC, all seven properties completed without a counterexample under
the small and medium profiles, yielding 140 completed runs. In
contrast, every measured execution under the large TLC profile reached
the predefined 1800-second timeout. This produced 70 censored
observations, corresponding to seven properties and ten measured
repetitions per property.

No out-of-memory or tool-error outcome was observed in the RQ1
measurements. The incomplete evidence is therefore concentrated
entirely in the large TLC profile.

\subsubsection{Answer to RQ1}
\label{subsubsec:rq1-answer}

Within the 350 completed bounded model-checking runs, no violation of
the seven evaluated properties was observed. This provides supporting
evidence for H1 within the completed configurations.

However, the 70 large-profile TLC measurements did not complete within
the predefined resource limit. Their absence of a counterexample
cannot be interpreted as successful property verification.

H1 is therefore partially supported under censoring. The result is
restricted to the completed bounded configurations and does not
constitute an unbounded proof of protocol correctness.

\subsection{RQ2: Mutation-based validation}
\label{subsec:rq2-results}

RQ2 evaluates whether the declared property suite detects predefined
scientific mutations that remove controls from the cross-shard
protocol.

The experiment contained ten scientific mutants: five instantiated in
the TLA+ model and five in the Alloy model. Each mutant was evaluated
through ten measured repetitions at the medium profile, producing 100
measured mutation runs.

All ten scientific mutants were detected through their designated
target properties, and detection remained consistent across all
measured repetitions. The resulting mutation score was

\begin{equation}
\mathrm{MutationScore}
=
\frac{10}{10}
=
1.0.
\label{eq:rq2-mutation-score}
\end{equation}

Table~\ref{tab:rq2-summary} summarizes the aggregate result.

\begin{table}[htbp]
\centering
\small
\caption{RQ2 scientific mutants, target properties, and detection results.}
\label{tab:rq2-summary}
\begin{tabular}{llllr}
\hline
Tool & Mutant & Target property & Detected & Runs \\
\hline
Alloy &
commit-after-abort &
TerminalStateIrreversibility &
Yes & 10/10 \\

Alloy &
credit-before-receipt &
DestinationCreditRequiresValidReceipt &
Yes & 10/10 \\

Alloy &
no-replay &
NoReceiptReplay &
Yes & 10/10 \\

Alloy &
quorum-bypass &
QuorumRequired &
Yes & 10/10 \\

Alloy &
timeout-without-release &
EventuallyReleasedAfterTimeout &
Yes & 10/10 \\

TLC &
commit-after-abort &
DecisionConsistency &
Yes & 10/10 \\

TLC &
credit-before-receipt &
DestinationCreditRequiresValidReceipt &
Yes & 10/10 \\

TLC &
no-replay &
NoReceiptReplay &
Yes & 10/10 \\

TLC &
quorum-bypass &
QuorumRequired &
Yes & 10/10 \\

TLC &
timeout-without-release &
EventuallyReleasedAfterTimeout &
Yes & 10/10 \\
\hline
\end{tabular}
\end{table}

The commit-after-abort defect is associated with different designated
target properties in the two formal representations:
\texttt{TerminalStateIrreversibility} in Alloy and
\texttt{DecisionConsistency} in TLA+. This difference follows the
tool-specific encodings of terminal consistency and does not affect
the mutant-level detection criterion.

The evaluated defect classes covered removal of replay protection,
destination credit without the required receipt condition, commit after
an abort-like terminal decision, timeout without release of source
funds, and quorum bypass.

Detection was evaluated at the mutant level rather than by counting
individual counterexample-producing repetitions as separate mutants.
Therefore, the 100 measured runs represent repeated evaluations of ten
scientific mutants, not 100 independent mutant definitions.

The repeated measurements also verified consistency of the logical
outcome: each mutant continued to expose its designated target-property
violation across all ten measured repetitions.

No absolute execution-time comparison between TLC and Alloy is used to
support RQ2. The two formal representations have different exploration
semantics, and the relevant result for this research question is
whether each predefined mutant exposes its designated property
violation.

\subsubsection{Answer to RQ2}
\label{subsubsec:rq2-answer}

All ten predefined scientific mutants were detected through their
target properties, yielding a mutation score of 1.0 with consistent
detection across the measured repetitions.

H2 is therefore supported for the predefined scientific mutant
catalogue.

This result demonstrates that the evaluated property suite is
sensitive to the injected defect classes considered in the experiment.
It does not establish completeness with respect to arbitrary
cross-shard protocol defects or mutations outside the predefined
catalogue.

\section{Discussion}
\label{sec:discussion}

\subsection{Bounded verification evidence}
\label{subsec:discussion-bounded-verification}

RQ1 provides bounded evidence that no declared property was violated in the
completed valid-model configurations. The property suite covers distinct
obligations, including replay protection, destination authorization, decision
consistency, value preservation, quorum enforcement, terminal-state
irreversibility, and release after timeout.

The evidential scope is determined by the explored bounds. Alloy completed all
three predefined profiles, whereas TLC-large was fully censored by the frozen
timeout policy. Those timeouts are not property violations; they delimit the
configurations for which bounded verification completed under the available
resource budget.

Accordingly, H1 supports only the absence of observed violations in completed
bounded configurations. It does not extend to unexplored states, censored
TLC-large executions, or an unbounded correctness claim.

TLC and Alloy provide complementary corroboration rather than interchangeable
proofs. Their different formal semantics and state representations prevent
agreement across completed profiles from implying identical explored spaces or
equivalent proofs.

\subsection{Mutation sensitivity and property adequacy}
\label{subsec:discussion-mutation}

RQ2 complements RQ1 by asking whether the declared properties react when
selected protocol controls are removed. All ten predefined scientific mutants
were detected by their designated target properties, producing a mutation score
of 1.0. This makes the valid-model results more informative because controlled
defects in replay protection, receipt-dependent credit, terminal decisions,
timeout release, and quorum enforcement trigger the expected properties.

Tool-specific encodings can express the same defect intention through different
properties. For example, commit after abort is associated with
\texttt{TerminalState\allowbreak Irreversibility} in Alloy and
\texttt{Decision\allowbreak Consistency} in TLA+. The mutant-level conclusion
therefore concerns targeted defect detection, not identical property
formulations across tools.

The perfect mutation score remains restricted to the predefined catalogue. It
does not establish completeness against arbitrary implementation defects,
modelling errors, or other cross-shard failure modes.

RQ1 and RQ2 thus support complementary claims: the declared properties hold in
the completed valid-model configurations, and the evaluated property suite is
sensitive to the selected protocol controls when they are deliberately
weakened.

\subsection{Implementation-model trace conformance}
\label{subsec:discussion-conformance}

RQ3 adds implementation-facing evidence by replaying observable Java traces
against the TLA+ abstraction. The valid catalogue was accepted and the
deliberately corrupted catalogue rejected as expected, connecting the formal
model with the implementation behaviors exercised by the frozen scenarios.

The negative controls strengthen this result because rejection was checked at
the expected diagnostic point rather than treated as a generic replay failure.
This provides evidence that deliberately introduced mismatches are localized at
the intended abstract-action boundary.

The result remains limited to the selected observable abstraction, scenario
catalogue, negative mutations, and frozen seeds. It therefore supports bounded
implementation-model trace conformance rather than a statement about every
possible Java execution.

In particular, this evidence must not be interpreted as a formal
refinement proof or as behavioral equivalence between the Java
implementation and the TLA+ specification. Establishing refinement
would require a stronger semantic relation covering the full relevant
transition systems, whereas the present experiment validates
correspondence for the predefined observable traces.

RQ3 therefore complements internal model checking and mutation analysis by
confronting the formal abstraction with observable implementation executions.

\subsection{Verification cost and scalability}
\label{subsec:discussion-cost}

RQ4 shows that bounded verification is constrained by computational cost as
well as logical outcome. Measurements increase within both formal tools, but
H4 remains unconfirmed in its preregistered nonlinear form.

For Alloy, all three ordered profiles completed, and median elapsed time and
memory increased monotonically from small to medium and large. The Spearman
association of 1.0 for both metrics is consistent with that ordering, but three
ordinal points cannot identify a functional growth law.

TLC exposes a sharper feasibility boundary. The small and medium profiles
completed with increases in elapsed time, memory, and explored state-space size,
whereas all seventy TLC-large measured runs reached the frozen 1800-second
timeout. Complete censoring precludes a completed-run median for the large
profile and a three-level characterization comparable to Alloy. The timeout is
therefore a valid censored observation, not an experimental failure, and
retaining it makes the evaluated verification boundary explicit.

Within the TLC medium-bound configurations, enabled fault profiles also changed
median time and memory across normal, insufficient-quorum, replay, and
timeout-oriented cases, indicating that verification cost depends on behavioral
structure as well as nominal bound size. These differences do not identify a
universal cost ordering, but they show that enabled behavior is part of the
verification envelope.

The evidence supports only a descriptive conclusion that bounded verification
cost increases for the evaluated configurations. It establishes neither a
nonlinear nor an asymptotic growth law. TLC and Alloy also use different
semantics, exploration mechanisms, and internal representations, so absolute
time, memory, or state-space measurements cannot support intrinsic
tool-superiority claims.

\subsection{Cross-layer evidence triangulation}
\label{subsec:discussion-triangulation}

The four research questions are intended to provide complementary evidence
rather than four independent demonstrations of protocol correctness. Each
addresses a different weakness in the validation argument.

RQ1 evaluates whether the declared properties survive the bounded
interleavings explored for valid formal models. Because absence of a
counterexample does not demonstrate sensitivity to relevant defects, RQ2
deliberately weakens selected protocol controls and checks whether their
designated properties detect the resulting violations.

Formal verification and mutation detection still concern formal
representations. RQ3 adds implementation-facing evidence by checking bounded
conformance of observable Java traces against the TLA+ abstraction, including
deliberately corrupted executions. This connects the formal layer with the
implementation without implying refinement or behavioral equivalence.

RQ4 makes the computational boundary explicit by characterizing verification
cost as the evaluated bounds become more demanding. The observed increase in
cost and complete TLC-large censoring show that the formal evidence is obtained
within a finite resource envelope.

Together, the layers form an evidence chain from bounded property verification
through mutation-based property validation and bounded implementation-model
trace conformance to verification-cost characterization. The chain is
complementary rather than cumulative proof: model checking remains bounded,
mutation analysis remains catalogue-dependent, trace replay remains limited to
the chosen abstraction and scenarios, and cost characterization remains tied to
the evaluated configurations. No layer removes the limitations of another.

The combined result is therefore not a general proof of cross-shard protocol
correctness. It is a reproducible multi-layer argument under one frozen
experimental protocol, connecting formal models, executable evidence, and the
resource conditions under which that evidence was obtained while preserving the
scope limits of each layer. This integration, rather than any individual
verification technique, is the methodological contribution emphasized here.


\section{Threats to Validity}
\label{sec:threats-validity}

\subsection{Construct validity}
\label{subsec:construct-validity}

A first construct-validity threat is the relationship between the TLA+ and
Alloy representations. Both address the same cross-shard protocol obligations
but use different state, transition, property, and analysis semantics.
Conceptually related properties are therefore not assumed to be syntactically
or semantically identical across tools. Claims are defined at the
protocol-obligation level, with mutation detection evaluated against the
designated property in each encoding. For example, commit after abort maps to
\texttt{TerminalStateIrreversibility} in Alloy and
\texttt{DecisionConsistency} in TLA+.

Verification-cost metrics introduce a second construct threat. Elapsed time and
maximum resident memory are observable for both tools, but state-space measures
are tool-specific. TLC distinct-state counts are not treated as equivalent to
Alloy's internal search representation. Cost trends are therefore interpreted
within each tool, and absolute timing, memory, or state counts are not used to
claim intrinsic tool superiority. Cross-tool observations are used only as
complementary evidence about the same protocol-level verification problem.

RQ3 introduces a third construct threat because it observes a deterministic
trace abstraction rather than the complete Java program state. The measured
construct is conformance of the selected observable abstraction, not semantic
equivalence between implementation and formal specification. The trace format,
scenario catalogue, mapping rules, and replay procedure were frozen before the
definitive campaign. Both valid and deliberately corrupted traces are tested,
and negative cases must match the expected rejection point rather than merely
fail replay.

Finally, mutation sensitivity is an operational proxy for property-suite
adequacy. The ten predefined scientific mutants exercise selected controls in
replay, receipt validation, terminal decisions, quorum enforcement, and timeout
handling, not an exhaustive defect taxonomy. A mutation score of 1.0 therefore
supports sensitivity to the predefined catalogue only and does not establish
complete defect detection or property-suite completeness. This distinction is
important because a perfect score on a targeted catalogue does not estimate
coverage of defect classes that were never represented experimentally.

\subsection{Internal validity}
\label{subsec:internal-validity}

Host load and execution conditions may affect verification time and memory.
The definitive campaign therefore used a dedicated native Linux environment,
with tasks executed serially at maximum parallelism one, reducing interference
between verification processes. This setup reduces, but cannot eliminate,
variation caused by operating-system scheduling and background activity.

Startup and runtime effects are addressed by separating warm-up executions from
measured repetitions. Warm-ups remain in the execution record but are excluded
from the summaries used to characterize measured verification cost.

Scenario generation and execution ordering could introduce nondeterminism.
Deterministic seeds, predefined configuration families, and the definitive
trace seed set were frozen before execution, and scenario definitions were not
changed in response to observed outcomes.

Measurement instrumentation is another possible confounder because
operating-system measurements and formal-tool parsers can fail independently
of the scientific execution. Instrumentation errors are therefore classified
separately, and reruns are permitted only under the predefined instrumentation
failure rule, not because a scientific outcome is unfavorable. This separation
prevents tool or instrumentation failures from being silently interpreted as
scientific results and prevents result-dependent reruns.

Timeout handling also affects internal validity. A run reaching the
1800-second limit has unknown eventual completion time and final verification
outcome. Timeout and out-of-memory results are retained as censored
observations, not converted to threshold-valued completions or successful
verification. The TLC-large timeouts therefore remain visible rather than
being selectively rerun.

More generally, the frozen protocol fixes tool versions, bounds, configuration
families, repetitions, seeds, execution order, resource limits, and treatment
of incomplete executions before interpretation of the definitive results,
reducing opportunities for result-dependent adaptation. It does not eliminate
environmental variation, but it makes the corresponding execution conditions
explicit and reproducible.

\subsection{External validity}
\label{subsec:external-validity}

The study evaluates one cross-shard commit protocol and its executable and
formal representations. Its findings characterize that protocol and
experimental implementation, not the full space of cross-shard mechanisms.
They should not be generalized directly to production blockchain systems with
different consensus protocols, network assumptions, cryptographic mechanisms,
execution environments, or adversarial models. The Java implementation is an
experimental simulator and validation artifact rather than a production
client.

The small, medium, and large formal profiles provide controlled bounded
configurations, not arbitrarily large deployments. Completion without a
counterexample within those profiles does not imply tractability or absence of
counterexamples under substantially larger bounds or different topologies. The
profiles vary shards, concurrent transfers, validators, quorum requirements,
receipt copies, and formal-analysis bounds only within the predefined campaign.

The mutation catalogue is likewise limited. Ten predefined scientific mutants
exercise selected protocol controls, while other defect classes and
interactions remain outside the evaluated catalogue. Mutation results therefore
generalize only to the evaluated defect classes and closely related protocol
obligations, not arbitrary bugs or adversarial behaviors.

Trace-conformance results are also scenario-bounded. The experiment uses ten
valid and ten deliberately corrupted scenario classes, each under thirty
deterministic seeds. This catalogue exercises multiple normal, timeout, replay,
quorum, concurrency, and corruption behaviors but is not an exhaustive sample
of Java executions. RQ3 therefore supports claims only for the evaluated trace
classes under the declared abstraction and bounds.

Finally, verification-cost observations depend on the selected models, tool
versions, bounds, failure configurations, execution environment, and resource
policy. They identify a scalability boundary for the evaluated setup, not
universal complexity laws for TLA+, Alloy, or cross-shard verification.
Accordingly, claims are restricted to the evaluated protocol, bounds,
scenarios, mutants, and tool configurations. The preserved protocol and
artifact support reproduction of those conditions, but do not expand the
population to which the results generalize.

\subsection{Conclusion validity}
\label{subsec:conclusion-validity}

RQ1 reports no counterexample among completed bounded verification executions,
but this is neither a statistical nor an unbounded proof of correctness. The
fully censored TLC-large profile prevents treating all scheduled RQ1 executions
as completed checks without violations. The RQ1 conclusion therefore applies
only to completed bounded configurations and retains censoring explicitly,
rather than treating incomplete executions as successful property checks.

Detection of all ten predefined scientific mutants yields a mutation score of
1.0 for the evaluated catalogue, but the mutants are not a random sample from
the universe of protocol or implementation defects. The score is descriptive
of sensitivity to the injected mutants, not an estimator of arbitrary future
defect-detection probability.

RQ3 likewise reports perfect observed classification for the twenty predefined
trace classes across thirty deterministic seeds per class. These seeds and
scenario classes are not a random sample of all possible Java executions.
Observed proportions and interval estimates therefore characterize the
evaluated experimental population rather than guaranteeing the same behavior
for arbitrary traces. The perfect observed classification rates therefore do
not constitute a population-level guarantee beyond the frozen catalogue.

RQ4 is limited by sparse and censored profile-level data. Alloy provides only
three completed ordered profiles. A Spearman rank association of 1.0 across
three levels establishes the observed monotonic ordering but cannot determine a
functional growth form or establish nonlinear scaling. TLC provides completed
measurements only for the small and medium profiles, while all measured
large-profile executions are censored by timeout, preventing a comparable
three-level analysis or estimation of a large-profile completion-time
distribution.

Censoring also affects descriptive summaries because medians computed from
completed executions characterize only those completions. If more demanding
runs are preferentially censored, completed-run summaries may underestimate the
computational burden of the full configuration. This is especially relevant
when censoring concentrates in the more demanding profiles, as observed for
TLC-large.

H4 is therefore interpreted as evidence of increasing verification cost
without sufficient evidence for the preregistered nonlinear-growth claim. No
asymptotic complexity law is inferred. More generally, conclusions emphasize
bounded property outcomes, mutation sensitivity, trace classification, effect
direction, and transparent censoring rather than significance claims not
supported by the frozen experimental design.


\section{Reproducibility and Artifact}
\label{sec:reproducibility-artifact}

\subsection{Artifact scope}
\label{subsec:artifact-scope}

The reproducibility artifact preserves the evidence path from the executable
protocol and formal models to the derived results reported in this paper,
rather than only source code or final tables.

The scientific-source layer contains the Java cross-shard implementation,
the TLA+ specifications and configurations used for bounded model checking,
and the Alloy models and commands used for bounded relational analysis.
These are the executable and formal objects evaluated by the campaign.

The frozen experimental-definition layer records the protocol version,
research questions and hypotheses, configuration families, bound profiles,
mutation cases, trace scenarios, deterministic seeds, repetitions, resource
limits, tool configurations, and rules for incomplete executions, separately
from the observed outcomes.

The evidence layer retains raw task results, derived datasets, and generated
tables and figures, preserving the link from reported summaries to the
underlying definitive observations.

The reproduction layer provides scripts and targets for structural checks,
isolated workspace preparation, representative verification and
trace-conformance smoke cases, analysis regeneration from preserved raw
results, and comparison with the reference artifact.

Integrity metadata, including manifests, content hashes, source-version
identifiers, and reproduction reports, connects the preserved campaign with
regenerated outputs and makes unexpected changes observable.

Together, the artifact supports inspection, representative partial
re-execution, analysis regeneration, and integrity verification. The
independent workflow does not require re-executing every definitive task and
distinguishes re-execution from regeneration based on preserved raw evidence.

\subsection{Reproduction workflow}
\label{subsec:reproduction-workflow}

The workflow separates structural validation, selective re-execution,
analysis regeneration, and integrity comparison.

First, it validates that the scientific sources, frozen experimental
definition, preserved raw results, derived outputs, and reproduction scripts
are available.

Second, it prepares an isolated workspace, records the source revision, and
creates regenerated-output directories without modifying the reference
campaign.

The third stage executes a representative smoke suite. The suite covers
a valid and a mutated TLC configuration, a valid and a mutated Alloy
configuration, one valid implementation trace, and one deliberately
corrupted trace. These executions exercise the principal formal and
implementation-facing paths of the artifact without repeating the full
definitive campaign.

The fourth stage regenerates the paper-level analysis from the preserved
raw experimental results. Derived datasets, tables, figures, and
research-question summaries are reconstructed using the same analysis
pipeline used for the reference artifact.

Fifth, regenerated outputs are compared with preserved references using
content hashes and reproduction gates to detect unexpected differences.

A reproduction verdict is issued only after structural checks, representative
executions, analysis regeneration, and integrity comparisons have completed.

Representative re-execution checks that the principal executable paths remain
operational, whereas analysis regeneration checks that paper-level evidence
can be reconstructed from the preserved definitive observations.

Neither operation is a new independent measurement of all definitive tasks.
Preserved raw results remain the source observations, while the smoke suite
only exercises selected artifact paths.

\subsection{Independent reproduction evidence}
\label{subsec:independent-reproduction}

The artifact was evaluated through an independent reproduction attempt using
a separate operating-system user, clone, and execution workspace from those
used to construct the reference campaign.

The reproduction used source revision
\texttt{6cd88c37\allowbreak 7afd23fe\allowbreak e4998882\allowbreak f91142d7\allowbreak 1e7d963e}
and packaged-artifact SHA-256 digest
\texttt{d464888e\allowbreak 9f3e5d8c\allowbreak c64ef5d2\allowbreak 2cc7b7c2\allowbreak 4f83e385\allowbreak 3f5825f1\allowbreak 8f23de26\allowbreak adf6a6e6}.

The reproduction workflow completed all ten predefined validation gates.
These gates covered artifact structure, workspace preparation,
representative formal and trace-conformance executions, analysis
regeneration, and integrity comparison. No reproduction incident was
recorded during the successful attempt.

All 32 expected output hashes matched their preserved references. The
successful reproduction therefore obtained 10/10 completed gates and 32/32
matching analysis artifacts, with zero recorded incidents.

The executable part of the reproduction consisted of the representative
smoke suite described above rather than a re-execution of the complete
definitive campaign. The full set of paper-level derived outputs was
instead regenerated from the preserved definitive raw observations and
checked against the reference artifact.

The attempt therefore demonstrates representative executable-path validation
and complete analysis regeneration from preserved raw evidence, not an
independent repetition of every definitive experimental measurement.

The reproduction provides process separation but not hardware independence:
a separate operating-system user, clone, and workspace were used on the same
native Linux host. This reduces dependence on the original working tree and
user-local execution state, while a second physical machine remains an
additional possible validation step.

\subsection{Integrity and traceability}
\label{subsec:artifact-integrity}

Source-version identifiers link the implementation, formal models, and
reproduction scripts to a repository state, while the frozen experimental
definition identifies the campaign associated with the reported results.

Preserved raw results are the primary empirical record. Derived datasets,
research-question summaries, tables, and figures are generated from these
observations through the documented analysis pipeline.

Integrity manifests record cryptographic hashes for analysis artifacts
expected to remain unchanged. Regenerated files are hashed again and compared
with the preserved references.

A matching hash establishes content identity for the corresponding
artifact. It does not independently establish scientific correctness,
but it provides a deterministic check that the reproduction pipeline
reconstructed the expected analysis output from the preserved evidence.

A hash mismatch is an observable reproduction difference requiring
investigation. The reproduction report records gate outcomes, integrity
comparisons, and incidents.

The resulting traceability chain can be summarized as:

\begin{quote}
source revision $\rightarrow$
frozen experimental protocol $\rightarrow$
preserved raw observations $\rightarrow$
derived analysis artifacts $\rightarrow$
integrity manifest $\rightarrow$
reproduction comparison.
\end{quote}

This chain supports bidirectional inspection, from paper-level values back to
preserved observations and from regenerated artifacts to the reference
manifest.

\subsection{Reproducibility boundaries}
\label{subsec:reproducibility-boundaries}

The successful reproduction was not a full experimental replication. It did
not re-execute every definitive task, and preserved raw observations remained
the empirical source for complete analysis regeneration.

It also does not establish hardware-independent reproducibility. The separate
user, clone, and workspace remained on the same native Linux host. A second
physical machine would provide a stronger test of environment portability.

Similarly, the artifact does not claim universal reproducibility across
arbitrary operating systems, processors, formal-tool versions, Java
runtimes, or resource configurations. Changes in these components may
affect execution behavior, verification cost, or generated outputs.

Finally, matching hashes establish content correspondence, not scientific
validity. The conclusions remain dependent on the experimental design,
formal abstractions, statistical treatment, and stated claim boundaries.


\section{Conclusions}
\label{sec:conclusions}

This study evaluated an executable cross-shard commit protocol through an
integrated, reproducible evidence workflow under one frozen experimental
protocol. Bounded property verification, mutation-based validation,
implementation-model trace conformance, verification-cost characterization,
and artifact reproduction provide complementary evidence rather than
cumulative proof, and no individual technique is claimed as novel in
isolation.

For RQ1, 350 of 420 measured bounded executions completed, while 70 TLC-large
executions were censored by timeout, with no property violations observed
among completed configurations. This is bounded evidence rather than an
unbounded correctness proof. For RQ2, all 10 predefined scientific mutants
were detected, giving a mutation score of 1.0 for the evaluated catalogue.
The result demonstrates sensitivity to those selected defect classes, not
completeness against arbitrary faults.

For RQ3, 600 trace executions produced the expected bounded conformance
outcomes: valid traces were accepted and deliberately corrupted traces were
rejected with the expected diagnostics. This does not establish refinement or
behavioral equivalence. For RQ4, completed Alloy profiles showed increasing
elapsed time and memory use, while all measured TLC-large executions reached
timeout. These observations characterize the evaluated verification envelope,
not a nonlinear or asymptotic scaling law or absolute TLC-versus-Alloy
superiority.

Taken together, these layers form a reproducible argument in which bounded
verification, targeted defect sensitivity, trace conformance, and cost
characterization address distinct limitations. Their combination strengthens
the evaluation without extending any layer beyond its stated evidential scope.

The analytical artifact was reproduced from a separate clone and workspace
under a separate operating-system user on the same native Linux host. The
attempt completed 10/10 validation gates and matched 32/32 expected
analysis-artifact hashes. This establishes process separation and deterministic
reconstruction under the documented environment, not hardware independence, a
complete rerun of all 1,272 scheduled tasks, or scientific correctness from
hash agreement alone.

Further work can extend the evaluated bounds, broaden the scientific mutant
catalogue and trace scenarios, and repeat reproduction on an independent
physical machine. These extensions would strengthen external validity while
preserving the bounded interpretation of the present evidence.


\bibliographystyle{elsarticle-num}
\bibliography{references}

\end{document}
